\scanpage{5}
1\textsuperscript{o}) \uline{$n$ pair} $=2\nu$, \uline{$n-1$ impair} \qquad $n-1=2\nu-1=2(\nu-1)+1$

\diagram{$\left.\begin{array}{ll}
 H^*(X_0):&1,\xi_0,\ldots,\xi_0^{\nu-1},[\xi_0^\nu,\lambda'_\nu],\frac12\xi_0^{\nu+1},\ldots\frac12\xi_0^{2\nu-1}\\[4pt]
 H^*(\bar X_1):&1,\xi_1\ldots,\xi_1^{\nu-1},\frac12\xi_1^\nu,\frac12\xi_1^{\nu+1},\ldots,\frac12\xi_1^{2\nu-1}
 \end{array}\right]\quad\begin{array}{l}\xi_0\leadsto\xi_1\\ \lambda'_\nu\leadsto0\end{array}$}

\[
\begin{aligned}
\text{Donc}\quad & H^i(X_0)\to H^i(\bar X_1)\quad\text{\uline{bijectif} \deleted{excepté} si }i\ne2\nu=n\\
& H^{2\nu}(X_0)\to H^{2\nu}(\bar X_1)\quad\text{\uline{surjectif}, noyau engendré par }\lambda'_\nu
\end{aligned}
\]

\[
\text{Donc}\quad\text{a)}\quad
\left\{\begin{aligned}
\Phi^{i}&=0\qquad \text{\deleted{sauf} si }i\ne(2\nu=n)\\
\Phi^{\unsure{2\nu|n}}&\simeq\mu(\unsure{-\nu})\qquad\text{(isom déterminé au signe près\ldots.)}
\end{aligned}\right.
\]

\[
\text{b)}\quad\left\{\begin{aligned}
&R^i f_*\quad\text{constant si }i\ne n\\
&R^n f_*\quad\text{a un sous-faisceau concentré en $0$, isom. à $\mu(\ldots)$.}\\
&\hspace{6em}\text{le quotient par ce dernier est constant}
\end{aligned}\right.
\]

\margin{c) $\pi$ opère trivialement sur $H^*(\bar X_1)$ et sur $\Phi^{2\nu}$}

2\textsuperscript{o}) $n$ impair $=2\nu+1$, \qquad $n-1=2\nu$

\diagram{$\left.\begin{array}{ll}
H^*(X_0):&1,\xi_0,\ \ ,\xi_0^{\nu-1},\xi_0^\nu,\frac12\xi_0^{\nu+1},\ldots\frac12\xi_0^{2\nu}\\[4pt]
H^*(\bar X_1):&1,\xi_1,\ldots\xi_1^{\nu-1},[\xi_1^\nu,\lambda_\nu],\frac12\unsure{\xi_1|\{_1}^{\nu+1},\ldots,\frac12\xi_1^{2\nu}
\end{array}\right]\quad\xi_0\leadsto\xi_1$}

\[
\begin{aligned}
\text{Donc}\quad & H^i(X_0)\to H^i(\bar X_1)\quad\text{bijectif si }i\ne2\nu=n-1\\
& H^{2\nu}(X_0)\to H^{2\nu}(\bar X_1)\quad\text{injectif, conoyau engendré par }\deleted{\unsure{\xi^\nu|\lambda'_\nu}}\inserted{\unsure{\ell_\nu|\lambda_\nu}}
\end{aligned}
\]

\[
\text{a)}\quad\left\{\begin{aligned}\Phi^i&=0\quad\text{si }i\ne2\nu+1=n\\
\Phi^n&\simeq\mu(?\ldots)\end{aligned}\right.
\]

\[
\text{b)}\quad\left\{\begin{aligned}
&R^i f_*\ \text{constant si }i\ne n-1\\
&R^{n-1}f_*\ \text{est ``pur'' \ldots.}
\end{aligned}\right.
\]

\margin{c) $\pi$ opère trivialement sur les $H^i(\bar X_1)$, pour $i\ne n-1$, mais non sur $\Phi^n$, \ldots.}

% archive sheet number: 4
